\begin{trapbox}
	\begin{description}[style=nextline, leftmargin=2.8em, labelsep=0.5em, itemsep=0.35em]
		\item[\textbf{\textcolor{CTrap}{易错点一（驻点误判）}}] 错误做法：认为 $f'(x_0)=0$ 则 $x_0$ 一定是极值点。
		\item[\textbf{\textcolor{CTrap}{正确理解（必要条件）}}] $f'(x_0)=0$ 仅是必要条件，还需 $f'(x)$ 在 $x_0$ 两侧变号。
		\item[\textbf{\textcolor{CTrap}{错因分析（主次混淆）}}] 混淆充分性与必要性，误将必要条件当作充分条件。
	\end{description}

	\medskip
	\begin{description}[style=nextline, leftmargin=2.8em, labelsep=0.5em, itemsep=0.35em]
		\item[\textbf{\textcolor{CTrap}{易错点二（同号忽略）}}] 错误做法：在检验极值点时，只计算 $f'(x_0)=0$，未检验两侧符号，直接认为是极值点。
		\item[\textbf{\textcolor{CTrap}{正确理解（两侧检验）}}] 必须检查导数在驻点两侧的符号是否改变。
		\item[\textbf{\textcolor{CTrap}{错因分析（缺乏验证）}}] 缺乏对充要条件的完整记忆，忽略变号步骤。
	\end{description}

	\medskip
	\begin{description}[style=nextline, leftmargin=2.8em, labelsep=0.5em, itemsep=0.35em]
		\item[\textbf{\textcolor{CTrap}{易错点三（计数错误）}}] 错误做法：认为极值点个数等于方程 $f'(x)=0$ 的实根个数。
		\item[\textbf{\textcolor{CTrap}{正确理解（变号计数）}}] 极值点个数等于变号根的个数。
		\item[\textbf{\textcolor{CTrap}{错因分析（遗漏条件）}}] 没有考虑不变号的根，直接对根计数。
	\end{description}
\end{trapbox}
